
\chapter{Modelo Quantum Inspired}
\label{app:quantum-inspired}

Este apéndice describe la arquitectura \emph{Quantum Inspired}, una red neuronal clásica que adopta un paradigma distinto al de matrices densas: utiliza una \textbf{factorización tensorial} inspirada en redes tensoriales usadas en física cuántica simulada clásicamente. La idea central es convertir el espacio de características embebidas en un tensor multidimensional y descomponerlo en una cadena de núcleos pequeños interconectados.

Esta estrategia es apropiada para sistemas con muchos grados de libertad (como la superficie de volatilidad) porque combina compresión estructural, regularización implícita y una lectura más trazable de cómo se propagan dependencias entre variables.

\section{Arquitectura general: flujo de datos desde entrada a predicción}

La arquitectura transforma un vector de entrada en una predicción escalar de volatilidad siguiendo tres bloques: preparación, núcleo Tensor-Train y agregación final. La Figura~\ref{fig:doc-models-families-quantum-inspired-1} resume el flujo.

\DiagramModelsFamiliesQuantumInspiredOne

\noindent La clave arquitectónica es que \textbf{el componente Quantum Inspired es la capa Tensor-Train}, no el conjunto completo. Las etapas anteriores preparan los datos; las posteriores agregan y refinan la salida. De este modo se separa la responsabilidad de cada bloque: preparación (embedding), transformación tensorial (TT) y agregación (post-TT dense).

\section{Etapas preparatorias y transformación estructural}

\subsection{Normalización e Embedding denso}

Sea $n_{\text{in}}$ el número de variables de entrada. La primera operación es una normalización por lote con transformación $n_{\text{in}} \rightarrow n_{\text{in}}$, que centra y escala cada mini-lote durante entrenamiento.

Después, el embedding denso aplica la transformación $n_{\text{in}} \rightarrow \prod_{k=1}^{d} n_k$, donde $d$ es el orden del tensor y $n_k$ es la dimensión del eje $k$ del tensor objetivo. Esta proyección es crítica porque:

\begin{itemize}
  \item \textbf{Captura relaciones no lineales:} modela interacciones entre variables desde el inicio.
  \item \textbf{Mejora y enriquecimiento:} la proyección denso-no-lineal reorganiza el espacio de características, destacando direcciones relevantes. Es análogo a un cambio de base adaptativo.
  \item \textbf{Evitación de cercanías espurias:} sin este paso, el reshape directo sería arbitrario. El embedding asegura que la interpretación de índices como ``adyacentes'' en el tensor tiene sentido matemático real.
\end{itemize}

El embedding se parametriza mediante:
\begin{equation}
\mathbf{z} = \sigma_{\text{emb}}\left(\mathbf{W}_{\text{emb}} \mathbf{x} + \mathbf{b}_{\text{emb}}\right),
\label{eq:qi-embedding}
\end{equation}
donde $\mathbf{x} \in \mathbb{R}^{n_{\text{in}}}$ es el vector de entrada y $\sigma_{\text{emb}}$ es la activación Swish. El resultado es $\mathbf{z} \in \mathbb{R}^{\prod n_k}$.

Se utiliza Swish en embedding porque suele mantener un flujo de gradientes más estable que alternativas más abruptas.

\subsection{Reshape a tensor de alta dimensionalidad}

El reshape realiza la transformación $\prod n_k \rightarrow n_1 \times n_2 \times \cdots \times n_d$. Es decir, el vector $\mathbf{z}$ se reinterpreta como un tensor de orden $d$:
\begin{equation}
\mathcal{Z} \in \mathbb{R}^{n_1 \times n_2 \times \cdots \times n_d},
\label{eq:qi-reshape}
\end{equation}
donde cada dimensión ($n_1, n_2, \ldots, n_d$) representa un grado de libertad abstracto que puede codificar factores de riesgo agrupados.

Este reshape no es un mero cambio de vista: es el pivote conceptual que permite pasar de un vector euclidiano a una estructura multiíndice. Una vez que $\mathbf{z}$ es visto como $\mathcal{Z}(i_1, i_2, \ldots, i_d)$ (con $1 \leq i_k \leq n_k$), se pueden aplicar contracciones tensoriales que de otro modo serían imposibles.

La importancia radica en que \textbf{genera un objeto de altísima dimensionalidad implícita}. Si por ejemplo \((n_1, n_2, n_3, n_4) = (2,2,2,2)\), entonces $\prod n_k = 16$ componentes se reinterpretan como un tensor de \(2^4 = 16\) entradas. La contracción TT que viene después explotará esta estructura multidimensional para comprimir ese espacio de manera exponencial.

\section{El núcleo Quantum Inspired: Descomposición Tensor-Train}

La descomposición tensor-train (TT) es una forma estructurada de representar un tensor multidimensional mediante un producto de matrices pequeñas encadenadas. Para un tensor $\mathcal{Z}(i_1, \ldots, i_d)$, la descomposición TT escribe:

\begin{equation}
\mathcal{Z}(i_1,i_2,\ldots,i_d) \approx G_1[i_1] \cdot G_2[i_2] \cdot \ldots \cdot G_d[i_d],
\label{eq:qi-tt-decomp}
\end{equation}

donde cada $G_k$ es un tensor de orden 3 con forma $(r_{k-1}, n_k, r_k)$:
\begin{itemize}
  \item $n_k$ es la dimensión del eje $k$ del tensor (determinada por la forma de $\mathcal{Z}$).
  \item $r_{k-1}$ y $r_k$ son los \textbf{rangos tensoriales} internos que conectan el núcleo $k-1$ con el $k$.
  \item Los rangos en las fronteras son forzados a 1: $r_0 = r_d = 1$, asegurando que el producto final es un escalar (o en nuestro caso, una componente de salida).
\end{itemize}

Para dejar el esquema claro desde el inicio: \textbf{cada salida TT tiene su propia cadena de núcleos}. Si la capa TT produce $m_{\text{out}}$ salidas, existen $m_{\text{out}}$ cadenas independientes, todas aplicadas sobre el mismo tensor de entrada.

Cada núcleo $G_k$ puede pensarse como una célula local que captura cómo la información asociada al eje $k$ interactúa con sus vecinos a través de los rangos. El encadenamiento de núcleos es el \textbf{contact chain} (cadena de contacto), mecanismo por el cual la información y las dependencias se propagan de un extremo del tensor al otro.

\subsection{La cadena de contacto (Contact Chain) y propagación de dependencias}

Durante el entrenamiento, todos los núcleos $G_1, G_2, \ldots, G_d$ aprenden simultáneamente. Los rangos internos $r_1, r_2, \ldots, r_{d-1}$ actúan como conexiones que permiten que gradientes y actualizaciones de parámetros se propaguen:

\begin{itemize}
  \item Un cambio en el núcleo $G_k$ afecta directamente a $r_{k-1}$ y $r_k$.
  \item Estos cambios, a su vez, repercuten en $G_{k-1}$ y $G_{k+1}$ a través de las contracciones tensoriales.
  \item De este modo, la red de dependencias es densa y altamente acoplada.
\end{itemize}

Esto contrasta radicalmente con una red densa donde cada capa aprende ``de forma más independiente'': aquí, cada núcleo no solo aprende su propia parametrización, sino que está en constante diálogo con sus vecinos a través de la estructura tensorial compartida.

\subsection{Implementación: estructura concreta de los núcleos}

En código, cada núcleo para un output dado es creado como:

\begin{minted}{python}
# Para cada salida out_idx in range(output_dim):
#   Para cada eje k del tensor de entrada:
#     Crear el núcleo con forma (r_k, n_k, r_{k+1})
tt_core = self.add_weight(
  name=f"tt_core_{out_idx}_{core_idx}",
    shape=(
    self.tt_ranks[core_idx],      # r_k
        mode_dim,                      # n_k  
    self.tt_ranks[core_idx + 1],  # r_{k+1}
    ),
    initializer=GlorotUniform(),  # Inicialización apropiada
    trainable=True,
)
\end{minted}

Sobre inicialización, en el código actual los \textbf{núcleos TT} se inicializan con \textbf{Glorot Uniform}. El hiperparámetro \texttt{kernel\_initializer} (\texttt{glorot\_uniform}/\texttt{he\_uniform}) se aplica a las capas densas de embedding y post-TT, no a los núcleos de \texttt{TensorTrainLayer}. En todos los casos, el objetivo es arrancar con una escala de pesos compatible con entrenamiento estable en las primeras épocas.

\subsection{Contracción: cómo se produce una salida TT}

Dado un tensor $\mathcal{Z}(i_1, \ldots, i_d)$ y una cadena de núcleos $\{G_1, \ldots, G_d\}$, la contracción TT produce una salida mediante:

\begin{equation}
\text{output} = \sum_{i_1=1}^{n_1} \cdots \sum_{i_d=1}^{n_d}
G_1[i_1] \cdot G_2[i_2] \cdot \ldots \cdot G_d[i_d] \cdot \mathcal{Z}(i_1, \ldots, i_d).
\label{eq:qi-contraction}
\end{equation}

Conceptualmente, la operación itera a través de los núcleos, contrayendo los índices de cada eje del tensor de entrada y propagando información a través de los rangos internos. Lo crucial es que esta contracción se realiza para \emph{cada} output: si hay $m_{\text{out}}$ salidas, existen $m_{\text{out}}$ cadenas independientes de núcleos, cada una contraída separadamente. Todas comparten la misma entrada $\mathcal{Z}$, pero aprenden distintos patrones de interacción.

En la ecuación~\ref{eq:qi-contraction}, $i_k$ recorre los índices del eje $k$, $G_k[i_k]$ es el bloque del núcleo $k$ asociado a ese índice, y el producto encadenado acumula la interacción entre ejes y rangos internos.

\subsection{Ejemplo concreto: relación entre dimensiones de entrada, rangos y salidas}

Supongamos:
\begin{itemize}
  \item Input original: $n_{\text{in}} = 5$ características (precio subyacente, strike, vencimiento, tasa, tipo).
  \item Tensor target: $(n_1, n_2, n_3, n_4) = (2, 2, 2, 2)$, así $\prod n_k = 16$.
  \item Embedding: capa densa $5 \to 16$.
  \item Reshape: los 16 componentes se organizan como tensor $2 \times 2 \times 2 \times 2$.
  \item Rangos TT: $(r_0, r_1, r_2, r_3, r_4) = (1, 2, 2, 2, 1)$.
  \item Salidas TT: $m_{\text{out}} = 8$ (especificado en configuración).
\end{itemize}

Aquí conviene distinguir dos magnitudes distintas: el valor $16 = \prod n_k$ describe el tamaño del \textbf{tensor de entrada} que recibe la capa TT, mientras que $m_{\text{out}} = 8$ describe el número de \textbf{cadenas TT independientes} y, por tanto, el número de características de salida producidas por la capa.

Para cada una de las 8 salidas:
\begin{itemize}
  \item Núcleo 1: forma $(1, 2, 2)$ — conecta ``antes'' (rango 0, siempre 1) con rango 1.
  \item Núcleo 2: forma $(2, 2, 2)$ — conecta rango 1 con rango 2.
  \item Núcleo 3: forma $(2, 2, 2)$ — conecta rango 2 con rango 3.
  \item Núcleo 4: forma $(2, 2, 1)$ — conecta rango 3 con ``después'' (rango 4, siempre 1).
\end{itemize}

Total de parámetros por cadena: $(1 \times 2 \times 2) + (2 \times 2 \times 2) + (2 \times 2 \times 2) + (2 \times 2 \times 1) = 4 + 8 + 8 + 4 = 24$ elementos.

Como la implementación crea \textbf{una cadena completa por cada salida} (no una única TT-matrix factorizando simultáneamente entrada y salida), el total en la capa TT es $8 \times 24 = 192$ parámetros. En el embedding, si contamos pesos y sesgo, son $5 \times 16 + 16 = 96$ parámetros. Esto sigue siendo compacto frente a una alternativa densa equivalente.

Se utiliza \textbf{Swish} como activación post-TT y en el embedding por su suavidad y por mantener un entrenamiento estable.

\section{Etapas finales: Agregación y salida}

Tras la contracción TT, tenemos $m_{\text{out}}$ características aprendidas. Estas pasan por una capa densa intermedia ($m_{\text{out}} \rightarrow u_{\text{dense}}$), seguida opcionalmente de Dropout, y finalmente por una capa de salida ($u_{\text{dense}} \rightarrow 1$) que genera la predicción escalar de volatilidad.

Esta última etapa ejecuta:
\begin{equation}
\hat{\sigma}(\mathbf{x}) = \mathbf{W}_{\text{out}} \cdot \phi(\mathbf{W}_{\text{post}} \cdot \mathbf{c}_{\text{TT}} + \mathbf{b}_{\text{post}}) + b_{\text{out}},
\label{eq:qi-final-output}
\end{equation}
donde $\mathbf{c}_{\text{TT}}$ son las $m_{\text{out}}$ salidas de la contracción TT.

Lectura rápida de los términos: $\mathbf{W}_{\text{post}}$ y $\mathbf{b}_{\text{post}}$ definen la proyección densa intermedia, $\phi(\cdot)$ representa la función de activación de ese bloque (en la configuración base se usa Swish), y $\mathbf{W}_{\text{out}}, b_{\text{out}}$ realizan la proyección final al escalar $\hat{\sigma}(\mathbf{x})$.

El Dropout es una regularización estocástica que desactiva aleatoriamente algunas neuronas durante entrenamiento, reduciendo co-adaptación y mejorando generalización.

\section{Justificación del enfoque Quantum Inspired}

La etiqueta ``quantum inspired'' se refiere a que adoptamos paradigmas de redes tensoriales originarios de la simulación clásica de sistemas cuánticos. En física cuántica, un estado multiqubit se describe mediante un tensor exponencialmente grande (en principio, $2^N$ componentes para $N$ qubits). Para computación tratable se usan descomposiciones tensoriales como tensor-train (conocida como \emph{matrix product state}) que comprimen esa representación en productos de matrices pequeñas. Esta misma intuición de representación compacta y aprendizaje supervisado con redes tensoriales se ha formalizado en literatura de referencia para ML clásico \citep{stoudenmire2016supervised,novikov2015tensorizing}.

La conexión con lo ``quantum'' está en la \textbf{estructura matemática} (red tensorial tipo MPS/TT), no en el hardware ni en el algoritmo de entrenamiento: aquí no hay qubits ni circuitos cuánticos, sino optimización clásica por gradiente sobre una parametrización inspirada en esas descomposiciones.

Aquí hacemos lo análogo en contexto clásico (predicción de volatilidad): reconocemos que aunque los datos sugieren un espacio de altísima dimensionalidad, es potencialmente compresible mediante factorización. Los grados de libertad del sistema (strike, vencimiento, subyacente, tasas) no son independientes; están acoplados por dinámicas de mercado. Capturar esa estructura acoplada mediante núcleos tensoriales es más eficiente que parametrización densa.

El contact chain (cadena de contacto) permite propagación bidireccional de dependencias durante entrenamiento: el modelo aprende no solo factores locales, sino cómo se \textbf{coordinan} esos factores en la cadena entera. Esto es esencial para capturar fenómenos como sonrisas de volatilidad.

La inferencia y entrenamiento son completamente clásicos (CPUs/GPUs estándar). No hay computación cuántica real: es puramente una arquitectura neuronal inspirada en estructuras tensoriales de física cuántica, implementada como red clásica bajo el mismo protocolo que otros modelos.

\section{Ventajas frente a una red densa estándar}

Una red secuencial convencional con capas densas tendría el siguiente flujo:

\begin{itemize}
  \item Input (5 dimensiones) $\to$ Dense (32 unidades) $\to$ Dense (24 unidades) $\to$ Output (1 dimensión).
\end{itemize}

Cálculo de parámetros (incluyendo bias):
$(5 \times 32 + 32) + (32 \times 24 + 24) + (24 \times 1 + 1) = 192 + 792 + 25 = 1009$ parámetros.

La arquitectura Quantum Inspired, con la misma aproximación de complejidad computacional:

\begin{itemize}
  \item Embedding: $5 \to 16$ ($5 \times 16 + 16 = 96$ parámetros).
  \item Tensor-Train (8 cadenas de salida, cada una con estructura $(2,2,2,2)$ y rangos $(1,2,2,2,1)$): $8 \times 24 = 192$ parámetros.
  \item Post-TT Dense: $8 \to 32 \to 1$: $(8 \times 32 + 32) + (32 \times 1 + 1) = 288 + 33 = 321$ parámetros.
  \item Total (sin BatchNorm): $96 + 192 + 321 = 609$ parámetros.
\end{itemize}

Con configuraciones similares de capacidad, el modelo TT es más compacto. Pero la ventaja no es solo la compresión de parámetros: es la \textbf{estructura inductiva} que impone la factorización tensorial. Mientras que una red densa puede representar cualquier función (dada suficiente capacidad), también puede aprender cualquier ruido. La factorización TT, al obligar a factorizar el espacio latente, \textbf{induce un sesgo inductivo} hacia soluciones que exhiben estructura de rango bajo, lo cual es a menudo una propiedad de sistemas complejos.

Adicionalmente, la propagación de gradientes en una arquitectura TT es más directa: los rangos internos actúan como canales de comunicación entre núcleos, reduciendo la necesidad de capas densas muy anchas para capturar interacciones complejas.
